% Blueprint: the Gram-point averaging lemma for Turing's method.
%
% References: Titchmarsh, *The theory of the Riemann zeta-function* (2nd ed.),
% \S 9.9 \cite{TitchmarshZeta}; Trudgian, *Improvements to Turing's method*
% \cite{TrudgianTuring}.

\chapter{The Gram-point averaging lemma}
\label{cha:gram_averaging}

Turing's method certifies the zero count $N(g_m)$ at a Gram point $g_m$ from the
integral bound of \Cref{cha:lehman_bound}. One of its analytic inputs is a bound on
the average, over a forward Gram window, of the normalised Riemann--Siegel phase
increment
\[
  x(t,l)\;:=\;\frac{\theta(t)-\theta(l)}{\pi},
\]
where $\theta=\texttt{riemannSiegelTheta}$ is the Riemann--Siegel theta of
\Cref{def:nf_theta}. This chapter proves that over a Gram interval this average is at
most $\tfrac12$. The bound is a consequence of the \emph{convexity} of $\theta$.

\begin{definition}[Normalised phase increment]
\label{def:ga_x}
\lean{gramCoord}
\leanok
For $l,t\in\mathbb{R}$ set
\[
  x(t,l)\;=\;\texttt{gramX}\,l\,t\;:=\;\frac{\theta(t)-\theta(l)}{\pi}.
\]
\end{definition}

\section{A chord bound for the integral of a convex function}
\label{sec:ga_chord}

The convexity input enters only through the following elementary trapezoidal
inequality, which has no special-function content and follows from mathlib's convexity
API.

\begin{lemma}[Chord integral bound]
\label{lem:ga_chord}
\lean{ConvexOn.integral_sub_left_le}
\leanok
Let $a<b$ and let $f:\mathbb{R}\to\mathbb{R}$ be convex and continuous on $[a,b]$.
Then
\[
  \int_a^b\bigl(f(t)-f(a)\bigr)\,dt\;\le\;\frac{(b-a)\bigl(f(b)-f(a)\bigr)}{2}.
\]
\end{lemma}

\begin{proof}
\leanok
Fix $t\in[a,b]$ and write $t=\alpha a+\beta b$ with
$\alpha=\frac{b-t}{b-a}\ge0$, $\beta=\frac{t-a}{b-a}\ge0$, $\alpha+\beta=1$. The
division-free chord inequality of convexity (mathlib
\texttt{ConvexOn.secant\_mono\_aux1} with the triple $a<t<b$) gives
\[
  (b-a)\,f(t)\;\le\;(b-t)\,f(a)+(t-a)\,f(b),
\]
i.e.
\[
  f(t)-f(a)\;\le\;\frac{t-a}{b-a}\,\bigl(f(b)-f(a)\bigr).
\]
Both sides are continuous, hence integrable on $[a,b]$; integrating and using
$\int_a^b(t-a)\,dt=\tfrac12(b-a)^2$,
\[
  \int_a^b\bigl(f(t)-f(a)\bigr)\,dt
  \;\le\;\frac{f(b)-f(a)}{b-a}\int_a^b(t-a)\,dt
  \;=\;\frac{f(b)-f(a)}{b-a}\cdot\frac{(b-a)^2}{2}
  \;=\;\frac{(b-a)\bigl(f(b)-f(a)\bigr)}{2}. \qedhere
\]
\end{proof}

\section{The derivative of the Riemann--Siegel theta function}
\label{sec:ga_deriv}

\begin{lemma}[First derivative of $\theta$]
\label{lem:ga_theta_deriv}
\lean{hasDerivAt_riemannSiegelTheta}
\leanok
\uses{lem:hz_theta_log}
For every $T>0$,
\[
  \theta'(T)\;=\;\tfrac12\,\operatorname{Re}\psi\!\left(\tfrac14+i\tfrac T2\right)-\tfrac12\log\pi,
\]
where $\psi=\texttt{Complex.digamma}$ is the digamma function.
\end{lemma}

\begin{proof}
\leanok
\uses{lem:hz_theta_log}
By the bridge \Cref{lem:hz_theta_log},
$\theta(T)=\operatorname{Im}\operatorname{Log}\Gamma(w(T))-\tfrac T2\log\pi$ with
$w(T)=\tfrac14+i\tfrac T2$ and $\operatorname{Log}\Gamma=\texttt{logGammaComplex}$
the holomorphic branch of \Cref{sec:hz_loggamma}. On the right half-plane
$\operatorname{Log}\Gamma$ is a holomorphic primitive of $\psi$
(\texttt{hasDerivAt\_logGammaComplex}), and $w'(T)=\tfrac i2$, so by the chain rule
$\tfrac{d}{dT}\operatorname{Log}\Gamma(w(T))=\psi(w(T))\cdot\tfrac i2$. Writing
$\psi(w)=A+iB$, its imaginary part is $\operatorname{Im}(\tfrac i2(A+iB))=\tfrac A2
=\tfrac12\operatorname{Re}\psi(w)$. Subtracting the derivative $\tfrac12\log\pi$ of
the smooth term $\tfrac T2\log\pi$ gives the claim.
\end{proof}

\section{Convexity of the Riemann--Siegel theta function}
\label{sec:ga_convex}

The main ingredient is monotonicity of $\operatorname{Re}\psi$ on the vertical
line, equivalently the non-positivity of $\operatorname{Im}\psi'$ in the upper
half-plane, where $\psi'=\texttt{deriv Complex.digamma}$ is the trigamma function. We
prove it through the trigamma \emph{recurrence}. Mathlib defines \texttt{Complex.digamma}
as $\operatorname{logDeriv}\Gamma$, so we first prove its differentiability.

\begin{lemma}[Differentiability of $\psi$]
\label{lem:ga_digamma_diff}
\lean{Complex.differentiableAt_digamma}
\leanok
$\psi=\texttt{Complex.digamma}$ is complex-differentiable on $\{\operatorname{Re}>0\}$.
\end{lemma}

\begin{proof}
\leanok
$\psi=\Gamma'/\Gamma$; on $\{\operatorname{Re}>0\}$ the numerator $\Gamma'$ is analytic
(as the derivative of the analytic, pole-free $\Gamma$) and the denominator $\Gamma$ is
analytic and non-vanishing (\texttt{Gamma\_ne\_zero\_of\_re\_pos}), so the quotient is
differentiable.
\end{proof}

\begin{lemma}[Trigamma recurrence]
\label{lem:ga_trigamma_rec}
\lean{Complex.deriv_digamma_add_one}
\leanok
\uses{lem:ga_digamma_diff}
For $\operatorname{Re}s>0$, $\psi'(s+1)=\psi'(s)-s^{-2}$.
\end{lemma}

\begin{proof}
\leanok
\uses{lem:ga_digamma_diff}
Differentiate the digamma functional equation $\psi(s+1)=\psi(s)+s^{-1}$
(\texttt{Complex.digamma\_apply\_add\_one}), valid on the open set $\{\operatorname{Re}>0\}$:
both sides are differentiable (\Cref{lem:ga_digamma_diff}), so their derivatives agree,
giving $\psi'(s+1)=\psi'(s)+\tfrac{d}{ds}s^{-1}=\psi'(s)-s^{-2}$.
\end{proof}

\begin{lemma}[Sign of the quadratic term]
\label{lem:ga_im_inv_sq}
\lean{Complex.im_inv_sq_nonpos}
\leanok
For $w$ in the closed first quadrant ($0\le\operatorname{Re}w$, $0\le\operatorname{Im}w$),
$\operatorname{Im}(w^{-2})\le0$.
\end{lemma}

\begin{proof}
\leanok
$\operatorname{Im}(w^{-2})=-\operatorname{Im}(w^2)/\lVert w^2\rVert^2
=-2(\operatorname{Re}w)(\operatorname{Im}w)/\lVert w^2\rVert^2\le0$.
\end{proof}

\begin{lemma}[Logarithmic bound on $\psi$]
\label{lem:ga_digamma_bound}
\lean{Complex.norm_digamma_le}
\leanok
\uses{thm:gs_digamma_stirling}
For $\operatorname{Re}z>0$ and $\lVert z\rVert\ge1$,
$\lVert\psi(z)\rVert\le\log\lVert z\rVert+(\pi+3)$.
\end{lemma}

\begin{proof}
\leanok
\uses{thm:gs_digamma_stirling}
From the order-$1$ digamma Euler--Maclaurin expansion
$\psi(z)=\operatorname{Log}z-(2z)^{-1}-\tfrac{B_2}{2}z^{-2}-R^{\psi}_3(z)$
(\Cref{thm:gs_digamma_stirling}), the principal term satisfies
$\lVert\operatorname{Log}z\rVert\le\lvert\log\lVert z\rVert\rvert+\lvert\arg z\rvert
\le\log\lVert z\rVert+\pi$ (using $\lVert z\rVert\ge1$), while each of the three remaining
terms has norm $\le1$ for $\lVert z\rVert\ge1$ (the remainder via
\Cref{thm:gs_remainder_bound_explicit}). Summing gives the bound.
\end{proof}

\begin{lemma}[Decay of trigamma along the shift]
\label{lem:ga_trigamma_decay}
\lean{Complex.tendsto_deriv_digamma_add_nat}
\leanok
\uses{lem:ga_digamma_bound, lem:ga_digamma_diff}
For fixed $s$ with $\operatorname{Re}s>0$, $\psi'(s+n)\to0$ as $n\to\infty$.
\end{lemma}

\begin{proof}
\leanok
\uses{lem:ga_digamma_bound, lem:ga_digamma_diff}
On the disc $\lvert z-(s+n)\rvert\le\tfrac12\operatorname{Re}(s+n)$ every point has
$\operatorname{Re}z\ge\tfrac12\operatorname{Re}(s+n)>0$, so $\psi$ is holomorphic there
(\Cref{lem:ga_digamma_diff}) and, by \Cref{lem:ga_digamma_bound}, bounded by
$\log\lVert z\rVert+(\pi+3)\le\log\!\bigl(2\lVert s+n\rVert\bigr)+(\pi+3)$ on the bounding circle.
Cauchy's derivative estimate (\texttt{Complex.norm\_deriv\_le\_of\_forall\_mem\_sphere\_norm\_le})
then gives $\lvert\psi'(s+n)\rvert\le\bigl(\log(2\lVert s+n\rVert)+\pi+3\bigr)/\bigl(\tfrac12\operatorname{Re}(s+n)\bigr)$,
which is $\le(2\log n+2K)/n\to0$ (with $K=\log(2(\lVert s\rVert+1))+\pi+3$), the limit
$\log n/n\to0$ being \texttt{Real.isLittleO\_log\_id\_atTop}.
\end{proof}

\begin{lemma}[Non-positivity of $\operatorname{Im}\psi'$]
\label{lem:ga_im_trigamma_nonpos}
\lean{Complex.im_deriv_digamma_nonpos}
\leanok
\uses{lem:ga_trigamma_rec, lem:ga_im_inv_sq, lem:ga_trigamma_decay}
For $\operatorname{Re}s>0$ and $\operatorname{Im}s\ge0$, $\operatorname{Im}\psi'(s)\le0$.
\end{lemma}

\begin{proof}
\leanok
\uses{lem:ga_trigamma_rec, lem:ga_im_inv_sq, lem:ga_trigamma_decay}
Iterating the recurrence \Cref{lem:ga_trigamma_rec} gives, for every $n$,
$\psi'(s)=\psi'(s+n)+\sum_{k=0}^{n-1}(s+k)^{-2}$, hence
$\operatorname{Im}\psi'(s)=\operatorname{Im}\psi'(s+n)+\sum_{k=0}^{n-1}\operatorname{Im}(s+k)^{-2}$.
Each summand is $\le0$ by \Cref{lem:ga_im_inv_sq} (the points $s+k$ lie in the closed
first quadrant), so $\operatorname{Im}\psi'(s)\le\operatorname{Im}\psi'(s+n)$; letting
$n\to\infty$ and using $\psi'(s+n)\to0$ (\Cref{lem:ga_trigamma_decay}) gives
$\operatorname{Im}\psi'(s)\le0$.
\end{proof}

\begin{lemma}[Monotonicity of $\operatorname{Re}\psi$ on a vertical line]
\label{lem:ga_re_digamma_mono}
\lean{Complex.monotoneOn_re_digamma_quarter_vertical}
\leanok
\uses{lem:ga_digamma_diff, lem:ga_im_trigamma_nonpos}
The map $T\mapsto\operatorname{Re}\psi\!\left(\tfrac14+i\tfrac T2\right)$ is monotone
increasing on $[0,\infty)$.
\end{lemma}

\begin{proof}
\leanok
\uses{lem:ga_digamma_diff, lem:ga_im_trigamma_nonpos}
With $w(T)=\tfrac14+i\tfrac T2$ one has, by the chain rule and \Cref{lem:ga_digamma_diff},
$\tfrac{d}{dT}\operatorname{Re}\psi(w(T))=\operatorname{Re}\bigl(\psi'(w(T))\cdot\tfrac i2\bigr)
=-\tfrac12\operatorname{Im}\psi'\!\left(\tfrac14+i\tfrac T2\right)\ge0$ for $T\ge0$, since
$\operatorname{Re}w(T)=\tfrac14>0$ and $\operatorname{Im}w(T)=\tfrac T2\ge0$ put $w(T)$ in
the closed first quadrant (\Cref{lem:ga_im_trigamma_nonpos}). A function with
non-negative derivative on an interval is monotone increasing.
\end{proof}

\begin{lemma}[Convexity of $\theta$]
\label{lem:ga_theta_convex}
\lean{convexOn_riemannSiegelTheta}
\leanok
\uses{lem:ga_theta_deriv, lem:ga_re_digamma_mono}
$\theta$ is convex on $[0,\infty)$.
\end{lemma}

\begin{proof}
\leanok
\uses{lem:ga_theta_deriv, lem:ga_re_digamma_mono}
By \Cref{lem:ga_theta_deriv} the derivative
$\theta'(T)=\tfrac12\operatorname{Re}\psi(\tfrac14+i\tfrac T2)-\tfrac12\log\pi$ is,
up to the additive constant $-\tfrac12\log\pi$ and the factor $\tfrac12$, the
function of \Cref{lem:ga_re_digamma_mono}, hence monotone increasing on $[0,\infty)$.
A differentiable function on an interval with monotone increasing derivative is
convex (mathlib \texttt{MonotoneOn.convexOn\_of\_deriv}), giving the claim.
\end{proof}

\section{The averaging bound}
\label{sec:ga_main}

\begin{theorem}[Gram-point averaging]
\label{thm:ga_main}
\lean{integral_gramCoord_le}
\leanok
\uses{def:ga_x, lem:ga_chord, lem:ga_theta_convex}
Let $0\le l$ and $h>0$, and suppose $[l,l+h]$ is a Gram interval, i.e.
$\theta(l+h)-\theta(l)=\pi$. Then
\[
  \int_l^{l+h} x(t,l)\,dt\;\le\;\frac h2,
  \qquad\text{equivalently}\qquad
  \frac1h\int_l^{l+h} x(t,l)\,dt\;\le\;\frac12 .
\]
\end{theorem}

\begin{proof}
\leanok
\uses{lem:ga_chord, lem:ga_theta_convex}
By \Cref{lem:ga_theta_convex}, $\theta$ is convex (and continuous) on $[l,l+h]\subseteq[0,\infty)$.
\Cref{lem:ga_chord} with $f=\theta$, $a=l$, $b=l+h$ gives
\[
  \int_l^{l+h}\bigl(\theta(t)-\theta(l)\bigr)\,dt
  \;\le\;\frac{h\bigl(\theta(l+h)-\theta(l)\bigr)}{2}
  \;=\;\frac{h\pi}{2},
\]
using the Gram spacing $\theta(l+h)-\theta(l)=\pi$. Dividing by $\pi$ and unfolding
$x(t,l)=(\theta(t)-\theta(l))/\pi$ (\Cref{def:ga_x}) gives
$\int_l^{l+h} x(t,l)\,dt\le\tfrac h2$; dividing further by $h>0$ gives the average
form. The constant is thus $C=\tfrac12$, sharp in the linear limit.
\end{proof}

\begin{theorem}[Gram-block averaging]
\label{thm:ga_block}
\lean{sum_integral_gramCoord_le}
\leanok
\uses{def:ga_x, thm:ga_main}
Let $p_0<p_1<\dots<p_k$ be an increasing block of points with $0\le p_0$, across each of
which $\theta$ increases by exactly $\pi$ (a \emph{Gram block} of $k$ consecutive Gram
intervals, $\theta(p_{i+1})-\theta(p_i)=\pi$ for $0\le i<k$). Then
\[
  \sum_{i=0}^{k-1}\int_{p_i}^{p_{i+1}} x(t,p_i)\,dt\;\le\;\frac{p_k-p_0}{2}.
\]
Equivalently, the average over the block of the \emph{fractional} phase increment
$\{(\theta(t)-\theta(p_0))/\pi\}$ is at most $\tfrac12$, independently of the number $k$
of Gram intervals spanned.
\end{theorem}

\begin{proof}
\leanok
\uses{thm:ga_main}
Each summand is measured from the \emph{left} endpoint $p_i$ of its own Gram interval, so
$x(t,p_i)$ is exactly the fractional part of the global increment there, and restarts at
$0$ at every Gram point. Since $0\le p_0\le p_i$ (from the strict inequalities) and
$\theta(p_{i+1})-\theta(p_i)=\pi$, \Cref{thm:ga_main} applied on $[p_i,p_{i+1}]$ gives
$\int_{p_i}^{p_{i+1}} x(t,p_i)\,dt\le\tfrac12(p_{i+1}-p_i)$. Summing over $0\le i<k$ and
telescoping $\sum_{i=0}^{k-1}(p_{i+1}-p_i)=p_k-p_0$ (\lean{Finset.sum_range_sub}) yields
the claim.

\emph{Remark (why not $k/2$).} A single chord over the whole block bounds the
\emph{unnormalised} increment $\int_{p_0}^{p_k}(\theta(t)-\theta(p_0))/\pi\,dt$ only by
$k(p_k-p_0)/2$, which is too weak for Turing's method (it does not force
$|S(g_n)|<1$ for $k\ge2$). The correct block quantity is the fractional one above, whose
average stays $\le\tfrac12$; this is what makes the ``Lehman bound over width $<\tfrac12$''
hypothesis pin the integer $S(g_n)$ to $0$.
\end{proof}
